\scanpage{77}
% printed form
% letterhead
UNIVERSITE\hfill Montpellier, le 8 Décembre 1977\\
DES SCIENCES ET TECHNIQUES\\
DU LANGUEDOC\\
------------\\
Service de la Scolarité\\
JP/GJ\\
Diffusion : Liste A\\
N°7040

\uline{AVIS DE SOUTENANCE D'UNE THESE DE DOCTORAT DE 3° CYCLE}

\uline{PHYSIQUE DU SOLIDE - APPLICATION A L' ELECTRONIQUE}

\uline{Nom et prénom du candidat} : ENSUQUE Louis

\uline{Titre de la thèse} :\textquotedbl{}Interprétation nouvelle des mécanismes de conduction dans les matériaux amorphes et désordonnés : application au B14Si\textquotedbl{}.

\uline{Date, heure et lieu de soutenance prévus} :\\
12 DECEMBRE 1977 à 10heures 30\\
Salle de réunion du C.E.E.S.

\uline{Nom du Directeur de thèse} :\\
Monsieur le Professeur PISTOULET.

\uline{Désignation du laboratoire où la thèse est préparée} :\\
Laboratoire d'Electronique C.E.E.S.

\uline{Analyse de l'ouvrage} :\\
Nous rappelons les théories actuelles concernant les corps amorphes, puis nous montrons les difficultés rencontrées pour l'interprétation de certains phénomènes physiques. Nous présentons un nouveau modèle de conduction dans les corps désordonnés et amorphes, qui tient compte de l'existence d'un désordre de composition à moyenne distance. Dans ce modèle un corps est caractérisé par la connaissance de deux quantités : énergie d'ionisation du niveau donneur ou accepteur et paramètre de désordre. Nous appliquons ce modèle au composé bore-silicium destiné à la conversion thermo électrique.

\uline{Durée approximative des recherches ayant abouti à la thèse} : 2 ans.

\uline{Composition du jury proposé} :\\
MM. : B. PISTOULET, Président\\
H. ANDRILLAT\\
J.L. ROBERT\\
J.M. DUSSEAU

Le Secrétaire Général,\\
\handwritten{\unsure{E. Siau|E. Si}}\\
\uline{E. SIAU}

% archive sheet number: 76
